\documentclass[11pt]{article}
\usepackage[T1]{fontenc}
\usepackage{lmodern,amsmath,amssymb,amsthm,booktabs,microtype,xcolor}
\usepackage[margin=1in]{geometry}
\usepackage[colorlinks=true,allcolors=blue!50!black]{hyperref}
\newtheorem{theorem}{Theorem}
\newtheorem{lemma}[theorem]{Lemma}
\newtheorem{proposition}[theorem]{Proposition}
\newtheorem{corollary}[theorem]{Corollary}
\DeclareMathOperator{\vol}{vol}
\DeclareMathOperator{\conv}{conv}
\newcommand{\R}{\mathbb R}
\newcommand{\CC}{\mathcal C}
\newcommand{\RK}{\mathcal P}
\title{Beyond point-generated product--join constructions\protect\\for projection-body volume}
\author{Robin Gounder}
\date{Working research draft, 9 October 2026}
\hypersetup{pdftitle={Beyond point-generated product-join constructions for projection-body volume},pdfauthor={Robin Gounder}}
\raggedbottom
\begin{document}
\maketitle
\begin{abstract}
We give an explicit rational twelve-dimensional seed whose iterated
products and joins have projection-volume root limit strictly between
$2.861$ and $2.8611$. In combination with the prior bound
$\Gamma_{\CC}\le e^{1.049}<2.855$ for the class generated from a point by
products, joins, and affine equivalences, this proves a strict asymptotic
separation between that entire class and unrestricted convex polytopes.
An explicit fixed-block construction exceeds every member of $\CC_n$
for every $n\ge15500$. A simpler eight-dimensional seed already gives
a limit above $2.8571$. The geometric seed family, product--join calculus,
recursive scheme, and comparison upper bound are attributed to prior
work. The additional contribution is their use with specified rational
seeds to obtain the stronger lower bounds and class separation.
Independent exact geometry and recurrence certificates accompany the
proof; global optimality and bibliographic priority are not asserted.
\end{abstract}

\section{The separation problem and main result}
For a full-dimensional convex body $K\subset\R^n$, use the normalization
\[
 h_{\Pi K}(u)=\vol_{n-1}(K\mid u^\perp),\qquad
 R_n(K)=\frac{\vol_n(\Pi K)}{\vol_n(K)^{n-1}}.
\]
Here $u$ is a unit vector and $K\mid u^\perp$ is its orthogonal projection.
Write $g(n)=n^n/n!$ and $c_n=(n+1)g(n)=R_n(T_n)$, where $T_n$ is an
$n$-simplex. The ratio is invariant under invertible affine maps.
For standard background on support functions, area measures, and affine
convex geometry, see Schneider~\cite{schneider}.
We omit the dimension subscript when it is unambiguous.
The standard product and join are taken in independent coordinate spaces:
\[
 A*B=\conv\bigl((A,0,0)\cup(0,B,1)\bigr).
\]
Let $\CC$ consist of polytopes obtained from a zero-dimensional point by
finitely many Cartesian products, joins, and invertible affine maps on
their affine hulls. Rank-dropping affine maps are outside this definition.
Set $\CC_n=\{K\in\CC:\dim K=n\}$ and
\[
 \Gamma_{\CC}=\limsup_{n\to\infty}\sup_{K\in\CC_n}R_n(K)^{1/n}.
\]
The limit in fact exists by the spectral reduction in~\cite{mxym-spectral},
but the limsup definition suffices below.
Let $U_n$ be the supremum of $R_n(K)^{1/n}$ over all full-dimensional
convex bodies $K\subset\R^n$.

The unrestricted maximization problem has a longer history.
Brannen~\cite{brannen} formulated the simplex maximum conjecture and
disproved a different proposed upper bound for centrally symmetric bodies.
Lutwak, Yang, and Zhang~\cite{lyz} developed affine-invariant functionals
for polytopes and derived projection-volume estimates. Saroglou~\cite{saroglou-cones}
established sharp three-dimensional maxima within the classes of zonoids,
cones, and double cones. These results concern the same volume ratio,
but do not supply the whole-class comparison used here.

Feng, Hu, Liu, and Xu~\cite{feng} already disprove the simplex upper bound
in every dimension at least nine using polytopes with at most $n+2$
facets. OpenAI's product example~\cite{openai} provides a particularly
simple later counterexample. Subsequent work in the mxym repository
gives the exact product--join calculus and a recursive construction with
root limit between $2.8534$ and $2.8535$~\cite[Section 9]{mxym-calculus}.
Its mixed Bellman bound~\cite{mxym-upper} controls every finite expression
in $\CC$, not just that particular recursion.
The question addressed here is whether the point-generated closure itself
can attain the best unrestricted exponential rate.

\begin{theorem}\label{thm:main}
There is an explicitly specified sequence of rational polytopes
$P_n\subset\R^n$ such that
\[
 \lim_{n\to\infty}R_n(P_n)^{1/n}=\beta,
 \qquad 2.861<\beta<2.8611.
\]
Consequently, using the prior upper theorem~\cite{mxym-upper},
\[
 \Gamma_{\CC}<2.855<\beta
 \le\liminf_{n\to\infty}U_n.
\]
There is also an explicit fixed-block sequence $E_n$ satisfying
\[
 R_n(E_n)>\sup_{K\in\CC_n}R_n(K)\qquad\text{for every }n\ge15500.
\]
The number $15500$ is a sufficient threshold, not a sharp one.
\end{theorem}

Thus the advantage concerns an entire construction class. For example,
the asymptotic ratio of $R_n(P_n)$ to its class supremum has lower root
at least $2.861/2.855>1.0021$. This is an exponential separation;
it is not a claim of optimality among all convex bodies.
For context, Zhang's reverse polar projection inequality~\cite{zhang}
and the Blaschke--Santal\'o inequality~\cite{meyer-pajor} give
\[
 R_n(K)\le\frac{\kappa_n^2 n^n}{\binom{2n}{n}},\qquad
 \limsup_{n\to\infty}U_n\le\frac{\pi e}{2},
\]
where $\kappa_n$ is the volume of the Euclidean unit ball; see
\cite[Introduction, equation (3)]{saroglou} for the combined estimate.
The limiting constant follows
from Stirling's formula. A substantial unrestricted gap therefore remains.

\section{Facet invariants and inherited operation rules}
For a polytope $K$ of dimension $d$ and volume $V$, let $v_i$ be its
outward facet area-normal vectors and let $\zeta_i$ be the facet area times
its support number. Equivalently, $\langle v_i,x\rangle=\zeta_i$ on facet $i$.
Define the determinant sums
\[
 \RK(K)=\sum_{|J|=d}|\det(v_j:j\in J)|,\qquad
 \mathcal S(K)=\sum_{|J|=d+1}|\det((v_j,\zeta_j):j\in J)|.
\]
Cauchy's projection formula gives $\RK(K)=\vol_d(\Pi K)$: the projection
body is a translate of the zonotope $\sum_i[0,v_i]$.
The determinant expression for zonotope volume and its application to
projection bodies are classical; see~\cite[Section 1]{lyz}
and~\cite[Section 3]{saroglou}.
The lifted invariant of~\cite[Section 2]{mxym-calculus} is
\begin{equation}\label{eq:invariant}
 a(K)=\frac{\mathcal S(K)}{dV\RK(K)},\qquad
 H(K)=\frac1{a(K)},\qquad Q(K)=\frac{a(K)R_d(K)}{g(d)}.
\end{equation}
It is affine invariant and satisfies $1/(d+1)\le a(K)\le1$.
We use different symbols for the geometric tilt parameters below and
the invariant $a$.

For positive-dimensional factors $A,B$ of dimensions $r,s$, the exact
identities in~\cite[Sections 3--5]{mxym-calculus} are
\begin{align}
 R(A\times B)&=R(A)R(B),&
 a(A\times B)&=\frac{ra(A)+sa(B)}{r+s},\label{eq:product}\\
 R(A*B)&=\frac{g(r+s+1)}{g(r)g(s)}R(A)R(B)(a(A)+a(B)),&
 a(A*B)&=\frac{a(A)a(B)}{a(A)+a(B)}.\label{eq:join}
\end{align}
These rules apply to the nonsimplex seeds as well as to members of
$\CC$. They, and the balanced recursion used below, are imported prior
results rather than new operation formulas.

\section{An explicit rational seed}
Let $n=2m$ with $m\ge1$, choose $0<u_1<\cdots<u_m<1$, and put
\[
 b=(-u_m,\ldots,-u_1,u_1,\ldots,u_m),\qquad
 K(u)=\{x\ge0:\ \boldsymbol1\cdot x+b\cdot x\le1,
                         \ \boldsymbol1\cdot x-b\cdot x\le1\}.
\]
For rational $u$, this is a rational realization of the central two-facet
section family underlying~\cite[Section 4]{feng}. We do not claim the family itself.
The symmetry exchanging paired coordinates interchanges the two upper
facets; it is not central symmetry of $K(u)$.

\begin{proposition}\label{prop:seed}
The polytope $K(u)$ has $n+1+m^2$ vertices and $n+2$ facets. Put
\begin{equation}\label{eq:I}
 I(u)=\sum_{i=1}^m
 \frac{u_i^{2m-2}}{(1+u_i)\prod_{j\ne i}(u_i^2-u_j^2)},\qquad
 p(u)=n+1+\sum_{j=1}^m(2j-1)u_j,\qquad
 s(u)=1+\sum_{j=1}^m u_j.
\end{equation}
Then
\begin{equation}\label{eq:seedstate}
 n!\vol_n(K(u))=I(u),\qquad
 R_n(K(u))=g(n)I(u)p(u),\qquad a(K(u))=\frac{s(u)}{p(u)}.
\end{equation}
In particular $Q(K(u))=I(u)s(u)$.
\end{proposition}
\begin{proof}
The vertices are the origin, the $n$ points
$e_i/(1+|b_i|)$, and, for each pair $b_i=-v<0<b_j=w$,
\[
 \frac{w}{v+w}e_i+\frac{v}{v+w}e_j.
\]
An active-set argument proves completeness: at a nonzero vertex one or
both upper facets must be active, leaving at most one or two positive
coordinates, respectively. In the latter case the signs must be opposite.
Each coordinate hyperplane contains the origin and $n-1$ independent axis
vertices. Each upper hyperplane contains a point with all coordinates
positive and $b\cdot x$ of the required strict sign. Thus every listed
inequality defines a facet.

On $b\cdot x\ge0$, the upper constraint with $+b$ determines feasibility;
the other upper constraint is redundant there.
Change variables $y_i=(1+b_i)x_i$ and use the standard simplex with slack
coordinate $y_0=1-\sum_{i=1}^n y_i$. Its truncating coefficients are
$t_0=0$ and $t_i=b_i/(1+b_i)$. The Jacobian is
$\prod_i(1+b_i)^{-1}$ and
$t_i-t_j=(b_i-b_j)/((1+b_i)(1+b_j))$ for $i,j\ge1$.
The simplex truncation identity in Appendix~\ref{app:truncation}, with
these factors cancelled, gives
\[
 n!\vol_n(K(u)\cap\{b\cdot x\ge0\})
 =\sum_{b_i>0}\frac{b_i^{n-1}}{(1+b_i)\prod_{j\ne i}(b_i-b_j)}.
\]
The two halves have equal volume. Pairing $b_j=\pm u_j$ in the last
denominator gives~\eqref{eq:I}.

Write $V=\vol_n(K(u))$. Each upper facet vector is a positive multiple
of $\boldsymbol1\pm b$. Reflection makes the two multiples equal.
The divergence identity $\sum_i\zeta_i=nV$, with zero support on coordinate
facets, makes each multiple $nV/2$. Balance of area-normal vectors then
gives $-nVe_i$ on coordinate facet $i$. After division by $nV$, the
horizontal generators and their lifted support coordinates are
\[
 (-e_i,0),\qquad
 \left(\frac{\boldsymbol1+b}{2},\frac12\right),\quad
 \left(\frac{\boldsymbol1-b}{2},\frac12\right).
\]
The horizontal $n$-minor sum is
\[
 1+n+\frac12\sum_{i<j}|b_i-b_j|
 =n+1+\sum_{j=1}^m(2j-1)u_j=p(u).
\]
The lifted $(n+1)$-minor sum is $1+\frac12\sum_i|b_i|=s(u)$.
Restoring the factors $(nV)^n$ and $(nV)^{n+1}$ in the two sums and
using~\eqref{eq:invariant} proves~\eqref{eq:seedstate}.
\end{proof}

Our principal seed is the following body with 49 vertices and 14 facets:
\begin{equation}\label{eq:seed12}
 K_0=K\left(\frac{17}{100},\frac{24}{100},\frac{32}{100},
                    \frac{41}{100},\frac{53}{100},\frac{69}{100}\right).
\end{equation}
In integer form its upper inequalities are
\[
 \sum_i(100\pm B_i)x_i\le100,\qquad
 B=(-69,-53,-41,-32,-24,-17,17,24,32,41,53,69).
\]
Direct rational evaluation, independently checked from its geometry,
gives
\begin{align}
 I_0&=\frac{1031929453906486710787861599001}
               {2386952337826268810834539955652},\label{eq:exactI}\\
 p_0&=\frac{768}{25},\qquad s_0=\frac{84}{25},\qquad
 a_0=\frac7{64},\qquad R_0=g(12)I_0\frac{768}{25}.\label{eq:exactstate}
\end{align}
For comparison, the simpler seed
$K(1/10,2/10,3/10,4/10)$ in dimension eight has 25 vertices and
\[
 I_0=\frac{231953}{378378},\qquad p_0=14,\qquad s_0=2,\qquad a_0=\frac17.
\]
Its initial ratio $R_8/c_8=231953/243243$ is below one. Nevertheless,
the same recursive amplification already exceeds the asymptotic ceiling
for $\CC$. A seed's immediate projection-volume root therefore does not
by itself describe its value under repeated operation rules.

\section{Recursive amplification and exact tail control}
Starting from either specified seed, apply the prior scheme
of~\cite[Section 9]{mxym-calculus}:
\begin{equation}\label{eq:recursion}
 K_{j+1}=(K_j\times K_j)*(K_j\times K_j),\qquad
 d_j=\dim K_j,\quad R_j=R_{d_j}(K_j),\quad a_j=a(K_j).
\end{equation}
The inherited identities give
\begin{equation}\label{eq:state-recurrence}
 d_j=\frac{(3d_0+1)4^j-1}{3},\qquad a_j=\frac{a_0}{2^j},\qquad
 R_{j+1}=R_j^4D_j,\quad
 D_j=2a_j\frac{g(4d_j+1)}{g(2d_j)^2}.
\end{equation}

\begin{proposition}[Exact seed ranking at fixed dimension]\label{prop:ranking}
For any two full-dimensional polytope seeds $A,B$ of the same positive dimension $d$,
the recursion~\eqref{eq:recursion} has positive root limits
$\beta_A,\beta_B$ and
\begin{equation}\label{eq:ranking}
 \left(\frac{\beta_A}{\beta_B}\right)^{3d+1}
 =\frac{a(A)R_d(A)^3}{a(B)R_d(B)^3}.
\end{equation}
Thus maximizing the limiting rate in the fixed-dimensional family
$K(u)$ is exactly equivalent to maximizing the rational objective
\[
 J(u)=I(u)^3p(u)^2s(u).
\]
This ranking concerns the specified recursion, not every operation schedule.
\end{proposition}
\begin{proof}
The central-binomial bounds used below imply that $D_j$ lies between
two positive constants depending only on $d,a_0$. The logarithmic
recurrence therefore gives a finite root limit for every seed.
For equal-dimensional seeds their multiplier ratio is $a(A)/a(B)$
at every level, so exact iteration gives
\[
 \frac{R_j(A)}{R_j(B)}
 =\left(\frac{R_d(A)}{R_d(B)}\right)^{4^j}
  \left(\frac{a(A)}{a(B)}\right)^{(4^j-1)/3}.
\]
Take $d_j$th roots and use $4^j/d_j\to3/(3d+1)$ to prove~\eqref{eq:ranking}.
Finally $aR^3=g(d)^3I^3p^2s$ by~\eqref{eq:seedstate}.
The factor $g(d)^3$ is constant at fixed dimension.
\end{proof}

\begin{lemma}\label{lem:tail}
For both seeds and every $j\ge0$, one has $3<D_j<7$ and
$R_j^{1/d_j}<3$. The roots $R_j^{1/d_j}$ strictly increase to a limit
$\beta$ with
\begin{equation}\label{eq:tail-enclosure}
 (3R_j^3)^{1/(3d_j+1)}<\beta<(7R_j^3)^{1/(3d_j+1)}.
\end{equation}
\end{lemma}
\begin{proof}
We extend the elementary tail argument in~\cite[Lemma 9.1]{mxym-calculus}
to these seed states. For $q\ge1$,
\[
 \frac{4^q}{2\sqrt q}\le\binom{2q}{q}
 \le\frac{4^q}{\sqrt{3q+1}}.
\]
Both follow by induction from equality at $q=1$.
For the lower bound the squared induction comparison has difference
$(2q+1)^2-4q(q+1)=1$; for the upper bound it has difference
$4(q+1)^2(3q+1)-(2q+1)^2(3q+4)=q$.
Set $q=2d_j$ and note that
\begin{equation}\label{eq:binomial}
 \frac{g(2q+1)}{g(q)^2}
 =\frac{4^q}{\binom{2q}{q}}
     \left(1+\frac1{2q}\right)^{2q}.
\end{equation}
The last factor is strictly between $2$ and $e$. Thus
\[
 4a_j\sqrt{6d_j+1}<D_j<4ea_j\sqrt{2d_j}.
\]
The closed dimension formula implies
\[
 a_j^2(6d_j+1)=a_0^2(6d_0+2-4^{-j})\ge a_0^2(6d_0+1),
\]
and $a_j^2(2d_j)<a_0^2(2d_0+2/3)$.
The exact checks
\begin{equation}\label{eq:hypotheses}
 16a_0^2(6d_0+1)>9,\qquad
 121a_0^2(2d_0+2/3)<49,\qquad
 7R_0^3<3^{3d_0+1}
\end{equation}
hold for both seeds. The elementary bound $e<11/4$ proves $D_j<7$.
Iterating $R_{j+1}<7R_j^4$ gives
\[
 R_j^3\le\frac{(7R_0^3)^{4^j}}7
 <\frac{3^{(3d_0+1)4^j}}7
 =\frac{3^{3d_j+1}}7<3^{3d_j}.
\]
Since $D_j>3>R_j^{1/d_j}$, the recurrence proves strict root increase.
The upper bound by three proves convergence. Iteration of the logarithmic
recurrence and the dimension formula yield the absolutely convergent series
\[
 \log\beta=
 \frac{\log R_j+\sum_{t\ge0}4^{-t-1}\log D_{j+t}}{d_j+1/3}.
\]
Using $3<D_{j+t}<7$ and $\sum_{t\ge0}4^{-t-1}=1/3$ proves
the two strict bounds in~\eqref{eq:tail-enclosure}.
\end{proof}

\begin{proposition}\label{prop:endpoints}
For the twelve-dimensional seed, $2.861<\beta<2.8611$.
For the eight-dimensional seed, $2.8571<\beta<2.8573$.
\end{proposition}
\begin{proof}
The finite certificate stores the full exact rational state $R_5$ for
each seed. The checker reconstructs it using~\eqref{eq:binomial}, without
importing the producer's factorial formula. It verifies by integer
comparison
\[
 \ell^{3d_5+1}<3R_5^3,\qquad 7R_5^3<h^{3d_5+1},
\]
with the following rational endpoints:
\begin{center}
\begin{tabular}{rrrr}
\toprule
Seed dimension & $d_5$ & $\ell$ & $h$\\
\midrule
8  & 8533  & $28571/10000$ & $28573/10000$\\
12 & 12629 & $2861/1000$   & $28611/10000$\\
\bottomrule
\end{tabular}
\end{center}
Together with Lemma~\ref{lem:tail}, these finite inequalities certify
the limits. No floating-point logarithm or numerically truncated infinite
product is used in acceptance of the endpoints.
\end{proof}

For the first assertion of Theorem~\ref{thm:main}, take the twelve-dimensional
family. For $n\ge144$, choose the largest $j$ with $d_j\le\sqrt n$, put
$k=\lfloor n/d_j\rfloor$, $r=n-kd_j$, and set
\[
 P_n=K_j^k\times T_r.
\]
Here $K_j^k$ denotes a Cartesian power; the factor $T_0$ is omitted.
For $n<144$ use $P_n=T_n$. Then $j\to\infty$,
$0\le r<d_j\le\sqrt n$, and $kd_j/n\to1$. Product multiplicativity gives
\[
 \frac{\log R_n(P_n)}n
 =\frac{kd_j}{n}\frac{\log R_j}{d_j}+\frac{\log c_r}{n}.
\]
Set $c_0=1$. The bound $c_r\le(r+1)e^r$ makes the last term tend to zero,
so the required all-dimensional limit is $\beta$. This padding argument
is the one used for the prior recursive family in~\cite[Corollary 9.3]{mxym-calculus}.

\section{The imported upper bound and an effective separation}
The mixed Bellman theorem~\cite{mxym-upper} states, for every $K\in\CC$,
with $D=n+1$ and $H=1/a(K)$,
\begin{equation}\label{eq:prior-upper}
 \log Q(K)\le\frac{87}{2000}\left(D-\frac{H^2}{D}\right)
              +\frac{11}{2000}\left(D-\frac{H^4}{D^3}\right)
 \le\frac{49}{1000}D.
\end{equation}
This is an imported theorem, not a new inequality of this note.
Its pinned exact certificate was replayed successfully: 19,900 finite
state rectangles, 15,568 tail cells, and the analytic large-dimension
constants. The complete English argument was also inspected. Exact
replay checks its finite proof obligations; it does not replace review
of its analytic argument.

Since $H\le n+1$, $R=g(n)HQ$, and $g(n)\le e^n$,
\begin{equation}\label{eq:finite-upper}
 \sup_{K\in\CC_n}R_n(K)
 \le e^{49/1000}(n+1)(e^{1049/1000})^n
 <\frac{1051}{1000}(n+1)\left(\frac{7137}{2500}\right)^n.
\end{equation}
The factorial bound follows from
$\sum_{i=1}^n\log i\ge\int_1^n\log t\,dt$.
The final bounds $e^{49/1000}<1.051$ and $e^{1049/1000}<2.8548$ are
certified by rational Taylor intervals.
Taking roots proves $\Gamma_{\CC}<2.855$, and hence the asymptotic
separation in Theorem~\ref{thm:main}.

\begin{proposition}\label{prop:effective}
Let $L=K_2$ for the seed~\eqref{eq:seed12}. Thus $\dim L=197$.
For $n=197k+r$, $0\le r\le196$, define $E_n=L^k\times T_r$.
Then $R_n(E_n)>\sup_{K\in\CC_n}R_n(K)$ for every $n\ge15500$.
\end{proposition}
\begin{proof}
Put $b_*=28581/10000=2.8581$. The exact recurrence certificate proves
$R_{197}(L)>b_*^{197}$.
The block has 224 facets and 46,118,408 vertices: under each recursive
step the facet count is multiplied by four and the vertex count changes
from $v$ to $2v^2$. The package supplies its sparse integer facet
inequalities, not its full vertex list.
The checker verifies all 197 rational residue inequalities
\[
 c_r>\frac{b_*^r}{3500}\qquad(0\le r\le196),\qquad c_0=1.
\]
It reconstructs the simplex values from $c_1=2$ and
$c_{r+1}/c_r=((r+2)/(r+1))(1+1/r)^r$, rather than the producer's
factorial formula. Consequently, for $k\ge1$,
\[
 R_n(E_n)>b_*^{n-r}c_r
 >\frac{b_*^n}{3500}.
\]
Comparison with~\eqref{eq:finite-upper} therefore reduces to
\[
 \left(\frac{9527}{9516}\right)^n>\frac{7357}{2}(n+1).
\]
The certificate checks this integer inequality at $n=15500$.
The ratio of consecutive left-to-right quotients is
$(9527/9516)(n+1)/(n+2)>1$ throughout $n\ge15500$.
This proves the assertion in every stated dimension.
\end{proof}

\section{Exact verification}
The accompanying standard-library Python package contains both exact
seeds, their vertex lists and oriented triangulations, the full rational
recurrence states, and separate geometric and growth checkers.
The producers use pulling triangulations; the standard geometric
framework is described in~\cite{triangulations}. The checkers validate
the resulting oriented chains independently of their construction.
For the twelve-dimensional seed, the geometric checker solves all 91
active systems, verifies 924 positively oriented simplices, cancels all
interior boundary faces, and verifies multiplicity one at a generic
rational interior point. The analogous counts for the eight-dimensional
seed are 45 active systems and 70 simplices.
All uncancelled faces lie on the body's boundary. Crossing an interior
codimension-one face therefore leaves the positive-chain multiplicity
unchanged, and a generic interior path avoids faces of codimension two.
Multiplicity is constant in the interior; its checked value one proves
coverage almost everywhere with no interior overlap.

The checker obtains volumes from simplex determinants and facet vectors
from boundary cofactors, then sums every horizontal and lifted maximal
minor. Thus it checks both $R_0$ and $a_0$ without assuming
Proposition~\ref{prop:seed}. The recurrence checker imports neither the
geometric producer nor its formula module and uses central binomial
coefficients instead of factorial ratios. It checks~\eqref{eq:hypotheses},
the rational limit endpoints, the dimension-197 block inequality, the
197 residue inequalities, and the dimension-15500 comparison. Explicit
exceptions keep the mathematical checks active under optimized Python.
The dimension-197 block also has an explicit inequality certificate.
A separate checker compares its 224 inequalities with a flattened
description using sixteen seed copies, four inner join coordinates,
and one outer join coordinate. It verifies a strict interior point.
This checks equivalence to the recursive expression; it does not
enumerate the block's vertices or independently recompute its volume.

Run \texttt{python verify\_all.py} for all new checks and nine unit tests.
The runner writes regenerated certificates and reports to temporary files
and leaves the delivered files unchanged. Empty or incomplete delivery
manifests and mathematically false comparison constants are rejected.
The command \texttt{python replay\_prior.py} fetches byte-pinned external
upper-bound inputs and replays their checker. This optional step needs network access
on first use. The delivery includes its successful replay report and the
input manifest. Formula and geometric data, unlike PDF build metadata,
regenerate byte-identically in the tested environments.

\section{Discussion}
The separation in Theorem~\ref{thm:main} identifies a limitation of the
point-generated class rather than a failure of products and joins
themselves. The same inherited operations amplify the new seeds beyond
the upper rate available to every finite expression in $\CC$.
Thus changing the initial geometric input can matter even when the
subsequent construction rules remain fixed. The result compares complete
construction classes; exceeding one previously known recursive sequence
alone would not establish this conclusion.

The eight-dimensional example illustrates why immediate seed performance
can be misleading. Its ratio $R_8/c_8$ is below one, yet its recursively
amplified rate exceeds the ceiling for $\CC$. The lifted invariant $a$
records information used by joins that the initial value of $R$ alone
does not capture. Proposition~\ref{prop:ranking} makes this observation
quantitative: at a fixed dimension and for the specified recursion, the
correct seed objective is $aR^3$, or equivalently $I^3p^2s$ in the stated
tilt family. This supplies an exact target for further seed search.
It does not justify comparing different dimensions by that objective or
claiming that the balanced recursion is the best possible schedule.

The asymptotic advantage is small at the level of a single root, but it
produces an exponential ratio of volumes as dimension grows. The fixed
block construction makes the comparison effective in every dimension
at least $15500$. This is a sufficient threshold obtained from a
particular block, residue bounds, and the imported class envelope; it
is not evidence that separation first occurs there. Nor does the lower
rate approach a known unrestricted optimum: the classical upper root
$\pi e/2$ leaves a substantial gap. The simplex-bound questions in
dimensions four through eight are unchanged.

Explicitness also needs a computational qualification. The block $L$
has a compact description by 224 facet inequalities but more than
46 million vertices. Its certified ratio is obtained from the checked
seed geometry and the inherited operation identities, rather than by
enumerating its vertices or constructing its projection body directly.
The package therefore makes these particular witnesses and their
comparison independently checkable; it does not provide an efficient
algorithm for optimizing projection-body volume over arbitrary polytopes.

The central geometric family is prior work~\cite{feng}, as are the lifted
invariant, operation rules, recursion, and tail method~\cite{mxym-calculus}.
The whole-class upper theorem~\cite{mxym-upper} is an essential external
dependency. Its exact finite obligations were replayed, but this does not
formalize its analytic argument or constitute external peer review.
The contribution relative to the inspected sources is the specified
rational seeds, their certified amplified rates, and the resulting
class separation and effective comparison. The supplementary literature
audit does not establish bibliographic priority.

Further progress could take several forms. A rigorous global search for
the objective $I^3p^2s$ would distinguish a good rational seed from an
optimum within a precisely stated family. Enlarging that family, changing
seed dimension, or allowing unbalanced schedules may improve the lower
rate. Sharper class envelopes and residue comparisons may lower the
effective dimension threshold. These are separate optimization problems;
none is resolved by the present certificates. External mathematical
review and proof-assistant formalization would strengthen confidence in
both the new argument and its imported dependencies.

\section{Conclusion}
An explicit rational twelve-dimensional seed, amplified by the prior
product--join recursion, gives a projection-volume root limit strictly
between $2.861$ and $2.8611$. Combined with the attributed upper theorem
for the point-generated class, this proves a strict exponential
separation from that entire class. A fixed 197-dimensional block and
simplex padding give a concrete comparison in every dimension
$n\ge15500$.

The examples show that the initial seed and its lifted invariant are
material to recursive growth. Exact geometry, recurrence certificates,
and a finite comparison make the claims reproducible without treating
numerical search as proof. The unrestricted extremal rate, optimal seed
and schedule, sharp separation threshold, and global bibliographic
priority remain open.

\section*{Generative AI disclosure}
OpenAI Codex agents contributed substantially to source discovery,
literature synthesis, and mathematical exposition.

\section*{Data and code availability}
The manuscript source, PDF, exact seed coordinates and triangulations,
dimension-197 facet inequalities, rational recurrence certificates,
verification programs, source audit, and reproduction instructions are
publicly available in the accompanying GitHub repository~\cite{data}.
The versioned release \texttt{v1.0.0} identifies the edition accompanying
this manuscript. Original Python code and generated certificate data are
provided under the MIT license; manuscript rights are stated separately.

From the release's \texttt{reproduction/} directory, run
\texttt{python -B verify\_all.py --regenerate} with Python 3.11 or newer.
All checks of the new constructions use only the standard library and
run without network access. The optional flag \texttt{--with-prior}
fetches and replays the byte-pinned external upper-theorem package;
its first use requires network access. The delivery manifest records
SHA-256 hashes, and verification leaves delivered files unchanged.
The release includes the successful external replay record. External
inputs remain attributed and are fetched from their pinned upstream
locations rather than redistributed as part of this package.

\appendix
\section{A simplex truncation identity}\label{app:truncation}
The following is a standard simplex section formula; see
Lasserre~\cite[Theorem 2.2]{lasserre}. We include a proof in the
normalization used here, using the Genocchi--Hermite representation of
divided differences~\cite[Section 9]{deboor}.
Let $t_0,\ldots,t_n$ be distinct real numbers, and let
$\Delta_n=\{y_i\ge0:\sum_{i=1}^n y_i\le1\}$, with
$y_0=1-\sum_{i=1}^n y_i$. Then
\begin{equation}\label{eq:truncation}
 \frac{\vol_n\{y\in\Delta_n:\sum_{i=0}^n t_i y_i\ge0\}}
      {\vol_n(\Delta_n)}
 =\sum_{t_i>0}\frac{t_i^n}{\prod_{j\ne i}(t_i-t_j)}.
\end{equation}
Here a zero node contributes zero.
For completeness, the elementary divided-difference integration identity is
\[
 \int_{\Delta_n} f^{(n)}\left(\sum_{i=0}^n t_i y_i\right)\,dy
 =\sum_{i=0}^n\frac{f(t_i)}{\prod_{j\ne i}(t_i-t_j)}.
\]
One proves it by induction using the fundamental theorem of calculus in
$y_n$. Multiplication by $t_n-t_0$ turns the integral into the difference
of the two $(n-1)$-simplex integrals on the faces $y_0=0$ and $y_n=0$.
The induction hypothesis gives exactly the divided-difference recurrence;
the initial case $n=0$ is evaluation at a point. Apply this identity to
$f(t)=t_+^n/n!$. Its $n$th derivative is the positive-half-line indicator
almost everywhere; the integration identity still holds by the same
fundamental-theorem argument, or by smooth approximation.
Since $\vol_n(\Delta_n)=1/n!$, equation~\eqref{eq:truncation} follows.

\begin{thebibliography}{99}
\bibitem{brannen}
N.~S.~Brannen, \emph{Volumes of projection bodies}, Mathematika
\textbf{43} (1996), 255--264.
\href{https://doi.org/10.1112/S002557930001175X}{Published article}.

\bibitem{deboor}
C.~de Boor, \emph{Divided differences}, Surveys in Approximation Theory
\textbf{1} (2005), 46--69.
\href{https://surveys-in-approximation-theory.com/papers/2/2.pdf}{Published full text}.

\bibitem{triangulations}
J.~A.~De Loera, J.~Rambau, and F.~Santos, \emph{Triangulations: Structures
for Algorithms and Applications}, Algorithms and Computation in
Mathematics 25, Springer, 2010.
\href{https://doi.org/10.1007/978-3-642-12971-1}{Publisher record}.

\bibitem{feng}
Y.~Feng, S.~Hu, W.~Liu, and L.~Xu, \emph{On the reverse projection
inequality}, author-deposited preprint (2026), MathSciDoc:2608.23002,
PDF dated 26 August 2026.
\href{https://archive.ymsc.tsinghua.edu.cn/pacm_paperurl/20260821091122099811781}{Primary archive record}.

\bibitem{data}
R.~Gounder, \emph{Beyond point-generated product--join constructions for
projection-body volume: manuscript, data, and code}, version 1.0.0 (2026).
\href{https://github.com/Kohze/projection-body-separation}{Reproduction repository};
\href{https://github.com/Kohze/projection-body-separation/releases/tag/v1.0.0}{Versioned release}.

\bibitem{lasserre}
J.~B.~Lasserre, \emph{Volume of slices and sections of the simplex in
closed form}, Optimization Letters \textbf{9} (2015), 1263--1269.
\href{https://doi.org/10.1007/s11590-015-0898-z}{Published article};
\href{https://arxiv.org/abs/1412.4617v1}{Versioned author preprint}.

\bibitem{lyz}
E.~Lutwak, D.~Yang, and G.~Zhang, \emph{A new affine invariant for
polytopes and Schneider's projection problem}, Transactions of the
American Mathematical Society \textbf{353} (2001), 1767--1779.
\href{https://doi.org/10.1090/S0002-9947-01-02726-X}{Published article};
\href{https://cims.nyu.edu/~yangd/papers/trans.pdf}{Author-hosted full text}.

\bibitem{meyer-pajor}
M.~Meyer and A.~Pajor, \emph{On the Blaschke--Santal\'o inequality},
Archiv der Mathematik \textbf{55} (1990), 82--93.
\href{https://doi.org/10.1007/BF01199119}{Published article}.

\bibitem{mxym-upper}
mxym (repository account), \emph{A certified mixed quadratic/quartic
envelope for projection-body growth}, research manuscript, dated
7 October 2026.
\href{https://github.com/mxym/math/blob/2fc3f25179cf1128c4bccbde65fd3a07c093ad6d/notes/mixed-bellman-product-join/proof.md}{Pinned proof};
\href{https://github.com/mxym/math/tree/2fc3f25179cf1128c4bccbde65fd3a07c093ad6d/notes/mixed-bellman-product-join}{Exact certificates and checker}.

\bibitem{mxym-spectral}
mxym (repository account), \emph{Asymptotic spectral reduction}, supplement
to the product--join calculus manuscript, research note (2026).
\href{https://github.com/mxym/math/blob/2fc3f25179cf1128c4bccbde65fd3a07c093ad6d/preprints/005-simplex-product-optimum/v2/ASYMPTOTIC_SPECTRAL_REDUCTION.md}{Pinned supplement}, Section 2.

\bibitem{mxym-calculus}
mxym (repository account), \emph{Projection-volume calculus for joins and
Cartesian products}, research manuscript, version 2, 7 October 2026.
\href{https://github.com/mxym/math/blob/2fc3f25179cf1128c4bccbde65fd3a07c093ad6d/preprints/005-simplex-product-optimum/v2/paper.md}{Pinned manuscript}, Sections 2--5 and 9.

\bibitem{openai}
OpenAI, \emph{A product counterexample to the simplex maximum for
projection-body volume}, research preprint, 24 September 2026.
\href{https://github.com/openai/math/tree/fd4aeeb2ee4fc729c18d98444fed42fd0529eeeb/preprints/A-product-counterexample-to-the-simplex-maximum-for-projection-body-volume-September-24-2026}{Pinned manuscript and source}.

\bibitem{saroglou-cones}
C.~Saroglou, \emph{Volumes of projection bodies of some classes of convex
bodies}, Mathematika \textbf{57} (2011), 329--353.
\href{https://doi.org/10.1112/S0025579311001860}{Published article}.

\bibitem{saroglou}
C.~Saroglou, \emph{On the shape of a convex body with respect to its
second projection body}, Advances in Applied Mathematics \textbf{67}
(2015), 55--74.
\href{https://doi.org/10.1016/j.aam.2015.03.004}{Published article};
\href{https://arxiv.org/abs/1409.4347v2}{Versioned author preprint}.

\bibitem{schneider}
R.~Schneider, \emph{Convex Bodies: The Brunn--Minkowski Theory}, second
expanded edition, Encyclopedia of Mathematics and its Applications 151,
Cambridge University Press, 2014.
\href{https://doi.org/10.1017/CBO9781139003858}{Publisher record}.

\bibitem{zhang}
G.~Zhang, \emph{Restricted chord projection and affine inequalities},
Geometriae Dedicata \textbf{39} (1991), 213--222.
\href{https://doi.org/10.1007/BF00182294}{Published article}.
\end{thebibliography}
\end{document}
